% paper1_results.tex -- Results of Paper 1, input by papers/paper1_indicators.tex.
% Every number is quoted from results/values.tex with \res{...}.
\section{Results}

The mean values of the six indicators are given in Table~\ref{tab:p1indicators}.
The integrated system had the highest soil organic carbon, water-stable
aggregates and microbial biomass carbon\index{microbial biomass carbon} at every site, and the conventional
system the lowest, while bulk density was lowest under the integrated system. Soil pH
differed little between systems and remained close to the starting value of
each site.

\input{\subfix{../tables/p1_indicators}}

The soil health index ranged from \res{p1-shi-min} to \res{p1-shi-max}
(Figure~\ref{fig:p1shi}). At the temperate site S1 it rose from
\res{p1-shi-S1-CONV} under the conventional system to \res{p1-shi-S1-CONS}
under the conservation system and \res{p1-shi-S1-INTG} under the integrated
system. The same order held at the semi-arid site S2, where the index was
lowest throughout (\res{p1-shi-S2-CONV}, \res{p1-shi-S2-CONS} and
\res{p1-shi-S2-INTG}), and at the humid tropical site S3
(\res{p1-shi-S3-CONV}, \res{p1-shi-S3-CONS} and \res{p1-shi-S3-INTG}). The
effect of the management system on the index was significant at each site
(S1: $F=\res{p1-F-S1}$, \res{p1-P-S1}; S2: $F=\res{p1-F-S2}$, \res{p1-P-S2};
S3: $F=\res{p1-F-S3}$, \res{p1-P-S3}; 2 and \res{p1-df-error} degrees of freedom).

\input{\subfix{../figures/fig1_shi}}
